\section{Conclusion}
\label{sec:conclusion}

This paper began with a real puzzle and a plausible prior. Ethereum's hourly
returns carry a marked downside tail asymmetry, the textbook leverage-effect
explanation is absent or inverted in crypto, and the field falls back on the
liquidation cascade of the leverage cycle, a mechanism that DeFi, alone
among markets, records transaction by transaction. A naive quantile local
projection appears to confirm the prior in full: a first-percentile response
\RatioTailMedImpact{} times the median on impact, deepening with the horizon,
robust across a battery of robustness checks, surviving multiple-testing correction.

A battery of diagnostics overturns that reading without overturning the
estimate. The coefficient is genuine in level. Both permutation nulls
leave mechanical overlap a minority share, with no signed bias. But its
apparent downside specificity does not survive. A symmetric, zero-skew placebo
reproduces the entire left--right tail gap at every horizon; the per-period
construction shows both tails loading equally; the headline ratio is an
artifact of an immobile median; and the whole signature replicates when
Bitcoin is substituted as the outcome. What the naive analysis detected is a
real but symmetric volatility channel, generic to crypto stress, dressed
up as downside fragility. This is the trap that symmetric heteroskedasticity sets
for quantile projections on overlapping cumulative returns, here sprung on
real data where the naive answer looked spectacular.

What survives is modest and worth having. Liquidations are a real,
symmetric predictor of ETH tail risk, in and out of sample. The effect is
incremental to
a nested quantile benchmark, not a replacement for a dedicated scale model;
the volatility
channel they drive confirms and extends the existing transaction-level
evidence to the conditional tails; and, net of volatility, the one-percent-tail
exceedance asymmetry (down minus up) is bounded below the smallest effect
size of interest under all three calibrations of the shock span, including the
strictest, so we can \emph{exclude} downside-specific amplification larger than
roughly \MdeEightyExc{} per unit of the log-liquidation shock. That is evidence of
absence at the economically relevant scale for the one-percent-tail
exceedance-asymmetry object under the stated calibrations of the shock span,
not a bare failure to reject. The small standardized residual at the one-percent tail is statistically
detectable but economically negligible, clearing two of the three calibrations
and borderline under the strictest, and it does not carry the conclusion.

The takeaway is bounded on every side. At the observable on-chain scale and
over this period, DeFi liquidations do not amplify ETH's downside tail
beyond their symmetric volatility effect; they nonetheless predict tail
risk, symmetrically; and a naive quantile local projection overreads that
symmetric volatility as directional fragility. This is not no effect: the
volatility channel is real. The short-horizon relation is bidirectional, so
the finding is not causal. It does not refute centralized-exchange cascades,
which our design cannot observe. And it does not claim that ETH's asymmetry
is fictional. That asymmetry is real; the observable DeFi channel simply adds no
downside-specific amplification to it, net of volatility.
The null is size-contingent: if the DeFi share of crypto leverage grows
materially, the channel could become aggregate-relevant, and the bound bears
re-examination as that share moves.

Four directions follow. We ran an unconditional-quantile (RIF) cross-check
(Section~\ref{sec:deconstruction}); because RIF-OLS is here algebraically tied
to the conditional model, a fully \emph{independent} unconditional
estimator would test the artifact from a separate angle. An identification
strategy
isolating exogenous variation in liquidation timing, whether an adjustable-rate
instrument or residualization at the quantile rather than the mean, would
license the causal statements this paper declines. And centralized
perpetual-futures data would let one test the spiral where it plausibly
lives, outside our perimeter. And broadening the on-chain liquidation panel
beyond the currently indexed lending protocols, to MakerDAO and DAI-vault
liquidations and to liquid-restaking collateral (Appendix~\ref{app:data}),
would close the two coverage gaps whose direction we can bound but not
eliminate here. The contribution we claim is correspondingly
sized: a reproducible on-chain liquidation panel, a portable methodological
caution that symmetric volatility masquerades as fragility in naive quantile
projections, a bounded null against a popular directional prediction, and a
symmetric tail-risk predictor, executed and stress-tested in full.
